\chapter{Basics}\label{chap:basics}

Some prose with \emph{emphasis}, \textbf{bold}, \texttt{code\_name} and a
citation~\cite{stacks}.  Maths is left untouched --- $\Spec A$, and
\[
  \poly{\Fq} \subset \Fq(T),
\]
where 100\% of the symbols survive.  See \url{https://stacks.math.columbia.edu}
and \href{https://leanprover-community.github.io}{the community site}.

\begin{quote}
  A blueprint is a plan, not a proof.
\end{quote}

\section{Rings}\label{sec:rings}

\begin{definition}[Polynomial ring]
\label{def:poly}
\lean{Example.polyRing}
\leanok
The polynomial ring over $\Fq$ is written $\poly{\Fq}$.
\end{definition}

\begin{lemma}
\label{lem:degree}
\uses{def:poly}
Every nonzero $f\in\poly{\Fq}$ has a degree.
\end{lemma}

\begin{lemma}
The zero polynomial has no degree.  This one carries no \texttt{label}, so it
gets a derived id.
\end{lemma}

\section{Notation}

\begin{tabular}{ll}
  \textbf{Symbol} & \textbf{Meaning} \\
  $\Fq$ & the field with $q$ elements \\
  $\poly{\Fq}$ & polynomials over it \\
\end{tabular}

\begin{description}
  \item[Ring] a commutative ring with unit.
  \item[Scheme] glued from rings.
\end{description}
